\pg{338} JII. Na ordem \emph{ontologica}, ou da \emph{realidade}, o ente \emph{necessario é anterior} ao \emph{contingente}, pois este é produzido por aquelle; mas na ordem \emph{logica}, ou do \emph{conhecimento}, o ente \emph{contingente é anterior} ao necessario, pois é por meio do \emph{contingente} que subimos ao conhecimento do \emph{necessario} \footnote{Na ordem \emph{ontologica}, \emph{o} ente \emph{mecessario} é anterior a0 contingente, pois aquelle é causa d'este, e toda a causa é anterior ao affeito. --- Na ordem logica, porém, dá-se o contrario. “O conceito do ente \emph{necessario} é posterior a0 conceito do \emph{contingente}. Porquanto, o nosso conhecimento começa pela percepção das coisas sensiveis, que estão mais ao nosso alcance. A nossa intelligencia, penetrando no ámago das coisas sensiveis, descobre nellas caracteres de mutabilidade, de defectibilidade, de potencialidade, e por isso de indifferenca para a existencia e para a não-existencia, e assim forma a idea do ente \emph{contingente}. E como o ente, indilferente para a existencia, não pode ser reduzido ao acto da existencia senão pela acção, immediata ou mediata, do ente, que possua por si, pela sua essencia, esse acto, assim da idea e pela, idea do ente \emph{contingente} nos elevamos a idea do ente \emph{necessario}.}.

%%>>> art-13

\section*{Artigo II.\\ Ente infinito e ente finito.}
\addcontentsline{toc}{section}{Artigo II. --- Ente infinito e ente finito.}

69, Ente infinito e ente finito. --- Ente \emph{infinito} é o que é

\emph{H} dotado de uma realidade ou perfeição illimitada, de modo que é a propria perfeição. --- Ente \emph{finito} é o que é dotado de uma realidade ou perfeigao \emph{limitada}, de modo que só possue a perfeição conveniente a um genero ou a uma especie. --- A idea do \emph{infi} nito é negativa-quanto-ao-termo;-mas- \emph{é-posiiva} quanto ao \emph{objecto}, pois representa um ente dotado de perfeição sem limite. A idea do \emph{finito}, mesmo considerada em relação ao \emph{objecto}, em parte é \emph{positiva}, emquanto exprime uma \emph{perfeição}, --- e em parte é \emph{negativa}, emquanto exprime uma perfeição \emph{limitada}, isto é, unida com a ausencia de ulterior perfeição (\textasciicircum{})\footnote{\emph{infinito} é o que não tem limites. Por isso, o \emph{infinito} não exclue toda e \emph{qualquer} outra realidade, mas só exclue tudo o que o pode limitar\pgn{339}na sua ordem. --- Um ente pode ser \emph{infinito} --- ou quanto á \emph{potencia}, que não é determinada por nenhum acto, como a \emph{materia prima}, --- ou quanto á \emph{forma}, cuja entidade não tem limite. No primeiro caso, o \emph{infinito} diz-se privativo, porque denota 'uma absoluta \emph{privação} ou ausencia de toda e qualquer perfeição num sujeito, que naturalmente devia ter uma perfeição; no segundo caso, o \emph{infinito} diz-se \emph{negativo}, porque exprime uma perfeição não coaretada por limites. --- O infinito negativo pode ser \emph{relativo} ou absoluto. É \emph{relativo}, quando é illimitado numa certa e determinada ordem de realidade ou perfeição; tal seria uma linha que na ordem do comprimento não tivesse fim. É absoluto, quando é illimitado em toda a ordem: de perfeição ou realidade, de modo que encerra todas as perfeições possiveis, e é a propria perfeição. É d'este infinito \emph{absoluto}, chamado simplesmente \emph{infinito}, que entendemos fallar neste artigo. --- D' onde se segue que tres elementos concorrem para a formação da idea do \emph{infinito}: um \emph{positivo}, que é a realidade ou perfeição, e dois \emph{negativos}, que são o \emph{limite} e a \emph{negação} do limite. Porquanto. o infinito, em si mesmo, é eminentemente positivo, porque é a propria perfeição: mas a vossa intelligencia não pode concebel-o senão de um inodo positivo-negativo, emquanto o concebemos como uma entidade positiva, da qual afastamos todo o limite, que é já uma negação, c concebemos a negação do limite, isto é, 2 negação da negação. E não podia ser d'outro modo. O infinito, na sua entidade positiva, escapa ao alcance da nossa intelligencia, e por isso, não podendo conceber o que elle é e não podendo exprimil-o pelo conceito, que formamos d'elle e que não é adequado, recorremos a negação, negamos os limites, pegamos a negação, e dizemos que não é \emph{limitado}, que é \emph{infinito}. Nominamos a Deus como o con- \textasciicircum{} cebemos; porque o termo, que denota um objecto, corresponde ao modo, por que o proprio objecto se concebe. Diz S. Thomaz: « Secundum quod aliquid a nobis cognosci potest, sic g\textasciicircum{}ríóbis potest nominari» (\emph{Sum}. \emph{Th}., p. L q. 13, a. 1). --- Do \emph{infinito} distiggue-se o \emph{indefinito}; porque este, embora possa augmentar e diminuir soB lim e por isso não tenha limites \emph{determinados}, comtudo na realidade é sempre finito.\par \emph{Finito é o} que tem limites, é o que não tem uma ulterior perfeição. Por isso, o \emph{finito}. como dizemos no texto, mesmo quanto ao \emph{objecto}, importa. --- uma \emph{affirmação}, porque \emph{effirma} ums \emph{perfeição}, --- e uma negação, porque \emph{nega} uma perfeição \emph{ulterior}. --- O ente \emph{finito é fa}! --- quanto á \emph{essencia}, \emph{que}, pertencendo a um certo genero ou a uma certa especie, não possue a perfeição dos outros generos ou das outras especies, --- e quanto á \emph{existencia}, a qual, por ser recebida numa \emph{potencia limitada}, \emph{como} é a essencia \emph{do} ente finito, não pode deixar de ser um acto limitado,}

63. A \textbf{nossa} intelligencia possue uma \textbf{idea} distincta do infinito. --- A nossa intelligencia possuirá uma idea \emph{distincta} do \emph{infinito}, se poder apontar o elemento, pelo qual o infinito se distingue dos seres finitos. Ora a nossa intelligencia aponta efectivamente o elemento, pelo qual o infinito se distingue dos seres finitos e que é a \emph{exclusão de todo o limite}.. Logo a nossa inteligencia possue uma idea \emph{distincta} do \emph{infinito}. --- Todavia essa idea, embora \emph{distincto}, \emph{é imperfeita} e \emph{inadequada}; porque \pg{339} 